Clutch performance in basketball refers to situations in which the outcome of the game is still uncertain and both teams 
are within 5 points difference in the last 5 minutes of the fourth quarter. This project analyses if the team’s clutch 
moments align with the overall teams quality or weather it is influenced by some statistical factors. Data for all 
30 NBA teams from the 2024–2025 regular season were collected directly from NBA.com.  A total of 3 datasets including 
clutch traditional statistics, clutch four-factor metrics, and regular-season traditional statistics were obtained and 
cleaned and preprocessed into a single dataset of 30 teams with 33 variables. The analysis used descriptive statistics 
and Pearson correlation to find relationships between clutch win percentage and key metrics such as shooting efficiency, 
turnovers, rebounding, and free-throw generation. Results show that regular season win percentage has a positive but 
imperfect relationship with clutch success. The strongest positive correlation with clutch win percentage was effective 
field goal percentage ($r \approx 0.51$), while turnovers and turnover percentage showed the strongest negative correlations
($\approx -0.50$). Overall, shooting efficiency and ball security appear to be the most important statistical factors associated 
with successful clutch performance.